\section{Diccionario de datos transformados}\label{sec:diccionario-transformado}
Esta sección documenta el conjunto de datos diario resultante de las
transformaciones de \cref{sec:transformaciones}. El archivo
\texttt{data/processed/sima\_transformado\_diario.parquet} constituye la entrada
de los análisis multivariados posteriores.

\subsection{Dimensión y cobertura}
El conjunto transformado contiene 26,691 observaciones día-estación y 34
columnas. Comprende 15 estaciones y cubre del 1 de enero de 2021 al 31 de diciembre
de 2025. Se transformaron y estandarizaron 16 predictores continuos medidos. Las
18 columnas restantes se conservaron en su escala original.

El criterio de \cref{sec:transformaciones} se reaplicó sobre la tabla diaria sin
modificar el método. De las 16 variables continuas, 13 requirieron Box-Cox y
estandarización z. \texttt{PRS\_media}, \texttt{RH\_media} y
\texttt{RH\_min\_diurna} cumplieron el criterio de sesgo desde la escala
original y solo se estandarizaron. Los valores faltantes se conservaron sin
imputación.

\subsection{Variables del conjunto diario}
La \cref{tab:dic-diario-continuas} presenta los 16 predictores continuos. Las
unidades corresponden a la escala original, antes de la transformación y la
estandarización.

\begin{table*}[t]
  \centering
  \caption{Predictores continuos transformados y estandarizados del conjunto
    diario. Las unidades corresponden a la escala original.}
  \label{tab:dic-diario-continuas}
  \footnotesize
  \setlength{\tabcolsep}{4pt}
  \begin{tabular}{lp{7.2cm}llr}
    \toprule
    Variable & Descripción & Unidad & Tratamiento & Faltantes \\
    \midrule
    \texttt{TOUT\_max} & Máximo diario de temperatura ambiente
      & $^{\circ}\mathrm{C}$ & Box-Cox + z & 1,561 \\
    \texttt{SR\_acum} & Acumulado diario de radiación solar
      & $\mathrm{kWh}/\mathrm{m}^{2}$ & Box-Cox + z & 1,123 \\
    \texttt{NO\_06\_09} & Media de NO entre las 06:00 y las 09:00 h
      & ppb & Box-Cox + z & 1,757 \\
    \texttt{NO2\_06\_09} & Media de NO$_2$ entre las 06:00 y las 09:00 h
      & ppb & Box-Cox + z & 1,776 \\
    \texttt{NOX\_06\_09} & Media de NO$_x$ entre las 06:00 y las 09:00 h
      & ppb & Box-Cox + z & 1,790 \\
    \texttt{CO\_06\_09} & Media de CO entre las 06:00 y las 09:00 h
      & ppm & Box-Cox + z & 1,937 \\
    \texttt{WSR\_media} & Media diaria de la velocidad del viento
      & km/h & Box-Cox + z & 1,008 \\
    \texttt{WSR\_min} & Mínimo diario de la velocidad del viento
      & km/h & Box-Cox + z & 1,008 \\
    \texttt{PRS\_media} & Media diaria de la presión atmosférica
      & mmHg & solo z & 886 \\
    \texttt{RH\_media} & Media diaria de la humedad relativa
      & \% & solo z & 2,649 \\
    \texttt{RH\_min\_diurna} & Mínimo diurno de la humedad relativa
      & \% & solo z & 2,575 \\
    \texttt{viento\_u\_media} & Media vectorial diaria del viento en la
      componente este-oeste & --- & Box-Cox + z & 1,447 \\
    \texttt{viento\_v\_media} & Media vectorial diaria del viento en la
      componente norte-sur & --- & Box-Cox + z & 1,447 \\
    \texttt{PM10\_media} & Media diaria de partículas PM$_{10}$
      & $\mu\mathrm{g}/\mathrm{m}^{3}$ & Box-Cox + z & 940 \\
    \texttt{PM2.5\_media} & Media diaria de partículas PM$_{2.5}$
      & $\mu\mathrm{g}/\mathrm{m}^{3}$ & Box-Cox + z & 5,965 \\
    \texttt{SO2\_media} & Media diaria de SO$_2$
      & ppb & Box-Cox + z & 1,933 \\
    \bottomrule
  \end{tabular}
\end{table*}

Con el 2025 completo el sesgo de \texttt{SR\_acum} sube a 0.53 y ya no
cumple el criterio, por lo que se transforma ($\lambda=0.755$); se conserva
como predictor continuo y no se convierte en binaria. Las 16 variables de la
\cref{tab:dic-diario-continuas} quedan dentro del criterio
$|\mathrm{sesgo}|<0.5$.

La \cref{tab:dic-diario-originales} reúne las 18 columnas que no se
transformaron. \texttt{msnm} se conserva porque representa un atributo de la
estación, con 15 valores repetidos, y no una medición diaria. \texttt{MDA8} y
\texttt{O3\_max\_1h} permanecen en ppb porque definen las variables objetivo y
sus umbrales deben conservar una interpretación física directa.

\begin{table*}[t]
  \centering
  \caption{Columnas conservadas en la escala original dentro del conjunto
    diario transformado.}
  \label{tab:dic-diario-originales}
  \footnotesize
  \setlength{\tabcolsep}{4pt}
  \begin{tabular}{llp{8.0cm}lr}
    \toprule
    Variable & Tipo & Descripción & Unidad & Faltantes \\
    \midrule
    \texttt{estacion} & nominal & Estación de monitoreo SIMA, con 15 niveles
      & --- & 0 \\
    \texttt{fecha} & fecha & Día calendario de la observación
      & --- & 0 \\
    \texttt{zona} & nominal & Zona geográfica del catálogo SIMA, con siete
      niveles & --- & 0 \\
    \texttt{msnm} & continua & Elevación de la estación sobre el nivel del mar
      & m & 0 \\
    \texttt{anio} & discreta & Año calendario
      & --- & 0 \\
    \texttt{mes} & discreta & Mes calendario
      & --- & 0 \\
    \texttt{temporada} & nominal & Temporada del año, con cuatro niveles
      & --- & 0 \\
    \texttt{tipo\_dia} & nominal & Día laborable, sábado, domingo o festivo
      & --- & 0 \\
    \texttt{festivo} & binaria & Indicadora de día de asueto
      & 0/1 & 0 \\
    \texttt{fin\_de\_semana} & binaria & Indicadora de sábado o domingo
      & 0/1 & 0 \\
    \texttt{llovio} & binaria & Indicadora de al menos una hora del día con
      precipitación & 0/1 & 800 \\
    \texttt{MDA8} & continua & Máximo diario de la media móvil de 8 h de O$_3$
      & ppb & 1,591 \\
    \texttt{ventanas\_validas} & discreta & Número de ventanas móviles de 8 h válidas
      durante el día, de un máximo de 24 & ventanas & 0 \\
    \texttt{O3\_max\_1h} & continua & Máximo horario de O$_3$ durante el día
      & ppb & 838 \\
    \texttt{excede\_anio\_1} & binaria & Excedencia de MDA8 sobre 65 ppb
      & 0/1 & 1,591 \\
    \texttt{excede\_anio\_3} & binaria & Excedencia de MDA8 sobre 60 ppb
      & 0/1 & 1,591 \\
    \texttt{excede\_anio\_5} & binaria & Excedencia de MDA8 sobre 51 ppb
      & 0/1 & 1,591 \\
    \texttt{excede\_1h} & binaria & Excedencia del máximo horario sobre 95 ppb
      & 0/1 & 0 \\
    \bottomrule
  \end{tabular}
\end{table*}

\subsection{Sesgo y curtosis}
La \cref{tab:dic-diario-sesgo} compara la forma de las distribuciones antes y
después del tratamiento. La estandarización z es una transformación lineal y no
modifica el sesgo ni la curtosis. Por ello, las tres variables que solo se
estandarizaron conservan los valores de la escala original.

\begin{table*}[t]
  \centering
  \caption{Sesgo y curtosis de exceso, con referencia normal igual a cero, de
    las variables continuas del conjunto diario antes y después del
    tratamiento.}
  \label{tab:dic-diario-sesgo}
  \small
  \setlength{\tabcolsep}{6pt}
  \begin{tabular}{lrrrrrr}
    \toprule
    & & & \multicolumn{2}{c}{Sesgo} & \multicolumn{2}{c}{Curtosis} \\
    \cmidrule(lr){4-5}\cmidrule(lr){6-7}
    Variable & Despl. & $\lambda$ & Antes & Después & Antes & Después \\
    \midrule
    \texttt{TOUT\_max} & 1.55 & 2.118 & $-0.86$ & $-0.16$ & 0.60 & $-0.49$ \\
    \texttt{SR\_acum} & 1.00 & 0.755 & 0.53 & 0.00 & 3.04 & 0.99 \\
    \texttt{NO\_06\_09} & 0.00 & $-0.103$ & 4.13 & 0.01 & 32.49 & $-0.09$ \\
    \texttt{NO2\_06\_09} & 1.00 & 0.279 & 1.67 & 0.02 & 6.89 & 0.40 \\
    \texttt{NOX\_06\_09} & 0.00 & 0.008 & 3.04 & 0.00 & 17.68 & 0.52 \\
    \texttt{CO\_06\_09} & 1.00 & $-0.320$ & 1.42 & 0.02 & 4.82 & $-0.53$ \\
    \texttt{WSR\_media} & 0.00 & 0.134 & 8.29 & 0.07 & 149.74 & 3.94 \\
    \texttt{WSR\_min} & 0.00 & 0.285 & 8.74 & 0.04 & 347.89 & 0.94 \\
    \texttt{PRS\_media} & --- & --- & 0.12 & 0.12 & $-0.16$ & $-0.16$ \\
    \texttt{RH\_media} & --- & --- & $-0.37$ & $-0.37$ & 0.19 & 0.19 \\
    \texttt{RH\_min\_diurna} & --- & --- & 0.45 & 0.45 & 0.15 & 0.15 \\
    \texttt{viento\_u\_media} & 116.16 & 2.955 & $-1.74$ & 0.29 & 27.19 & 5.23 \\
    \texttt{viento\_v\_media} & 99.96 & 2.956 & $-1.52$ & 0.20 & 27.68 & 4.86 \\
    \texttt{PM10\_media} & 0.00 & 0.165 & 1.58 & 0.01 & 4.98 & 0.48 \\
    \texttt{PM2.5\_media} & 0.00 & 0.021 & 2.04 & 0.00 & 9.86 & 0.03 \\
    \texttt{SO2\_media} & 0.00 & $-0.202$ & 3.63 & $-0.01$ & 36.70 & 0.47 \\
    \bottomrule
  \end{tabular}
\end{table*}

La agregación diaria redujo la curtosis de exceso de las componentes
\texttt{viento\_u\_media} y \texttt{viento\_v\_media} de 37.5 a 27.4 y de 38.7
a 28.3, respectivamente. Esta disminución equivale aproximadamente a 27 \% y
no elimina por sí sola las colas pesadas. Box-Cox las reduce a 5.25 y 4.97.

Las mayores curtosis antes de la transformación corresponden a
\texttt{WSR\_min}, con 342.92, y \texttt{WSR\_media}, con 142.35. Después de
Box-Cox disminuyen a 0.93 y 4.02. Aunque todas las variables cumplen el criterio
de sesgo, la curtosis residual del bloque de viento debe considerarse al
interpretar el análisis discriminante lineal.
